\documentclass[aspectratio=169,10pt]{beamer}
\makeatletter\def\input@path{{../../estilo/}{../estilo/}{estilo/}}\makeatother
\usepackage{devmobile}
\modulo[Kotlin]{02}{Kotlin Essencial}{Tipos, null safety, controle de fluxo, funções, orientação a objetos, coleções e funções de escopo}

\begin{document}
\capa

\begin{frame}{Objetivos do módulo}
  \begin{conceito}[Ao final deste módulo você será capaz de]
    \begin{itemize}
      \item Declarar variáveis e usar os tipos básicos e a inferência de tipos do Kotlin.
      \item Escrever código seguro contra \cod{NullPointerException} com \emph{null safety}.
      \item Usar \cod{if} e \cod{when} como expressões, além de \emph{ranges} e laços.
      \item Criar funções com parâmetros padrão, funções de extensão e lambdas.
      \item Modelar dados com \cod{data class}, \cod{enum}, \cod{sealed} e interfaces.
      \item Manipular coleções com operações funcionais e funções de escopo.
    \end{itemize}
  \end{conceito}
  \begin{dica}[Onde praticar]
    \url{https://play.kotlinlang.org} roda Kotlin no navegador, sem instalar nada. Os \emph{Kotlin Koans} trazem exercícios guiados.
  \end{dica}
\end{frame}

\sumario

% =============================================================================
%  Fundamentos compartilhados --- Módulo 02: linguagens de programação
%  Incluído pelas trilhas Kotlin e Flutter com % =============================================================================
%  Fundamentos compartilhados --- Módulo 02: linguagens de programação
%  Incluído pelas trilhas Kotlin e Flutter com % =============================================================================
%  Fundamentos compartilhados --- Módulo 02: linguagens de programação
%  Incluído pelas trilhas Kotlin e Flutter com \input{../fundamentos/f02-linguagens}
% =============================================================================
\section{Fundamentos: linguagens de programação}

\begin{frame}{Paradigmas de programação}
  \begin{columns}[T]
    \begin{column}{0.54\textwidth}
      \scriptsize
      \begin{tabularx}{\linewidth}{@{}l L@{}}
        \toprule
        \cabecalho{Paradigma} & \cabecalho{Ideia central} \\
        \midrule
        Imperativo & sequência de comandos que alteram o estado (\emph{faça isto, depois aquilo}) \\
        Orientado a objetos & dados e comportamentos agrupados em objetos que trocam mensagens \\
        Funcional & funções puras, valores imutáveis, composição de funções \\
        Declarativo & descreve \textbf{o que} se quer, não \textbf{como} obter (SQL, interfaces em Compose e Flutter) \\
        \bottomrule
      \end{tabularx}
    \end{column}
    \begin{column}{0.44\textwidth}
      \begin{definicao}[Paradigma]
        Estilo de organizar o raciocínio e o código. Linguagens modernas são \textbf{multiparadigma}: Kotlin e Dart misturam imperativo, objetos e funcional.
      \end{definicao}
      \begin{dica}
        Use o paradigma que deixa cada parte mais clara: objetos para modelar o domínio, funções para transformar dados, declarativo para descrever a interface.
      \end{dica}
    \end{column}
  \end{columns}
  \begin{historia}
    Fortran (1957) era imperativo; Lisp (1958), funcional; Simula (1967) e Smalltalk (anos 1970) criaram a orientação a objetos. Kotlin (2011, JetBrains) e Dart (2011, Google) herdaram ideias de todos.
  \end{historia}
\end{frame}

\begin{frame}{Sistemas de tipos}
  \begin{columns}[T]
    \begin{column}{0.5\textwidth}
      \begin{definicao}[Tipo]
        Conjunto de valores possíveis e das operações permitidas sobre eles. Somar dois \cod{Int} faz sentido; somar um \cod{Int} a uma \cod{Bancada}, não.
      \end{definicao}
      \scriptsize
      \begin{tabularx}{\linewidth}{@{}l L@{}}
        \toprule
        \cabecalho{Classificação} & \cabecalho{Significado} \\
        \midrule
        Estática & tipos verificados na compilação (Kotlin, Dart, Java) \\
        Dinâmica & tipos verificados na execução (Python, JavaScript) \\
        Forte & conversões perigosas são proibidas ou explícitas \\
        Inferência & o compilador deduz o tipo (\cod{val x = 3}) \\
        \bottomrule
      \end{tabularx}
    \end{column}
    \begin{column}{0.48\textwidth}
      \begin{conceito}[Por que tipagem estática ajuda]
        Muitos erros aparecem \textbf{antes} de o app rodar; a IDE completa código e refatora com segurança; o tipo documenta a intenção.
      \end{conceito}
      \begin{pordentro}[Solidez (\emph{soundness})]
        Um sistema de tipos é sólido quando garante que, se o programa compila, uma variável nunca terá em execução um valor de outro tipo. O Dart 3 é sólido, inclusive quanto a nulos.
      \end{pordentro}
    \end{column}
  \end{columns}
\end{frame}

\begin{frame}{O ``erro de um bilhão de dólares'': o nulo}
  \begin{columns}[T]
    \begin{column}{0.5\textwidth}
      \begin{historia}
        Tony Hoare introduziu a referência nula na linguagem ALGOL W em 1965. Décadas depois, chamou-a de seu ``erro de um bilhão de dólares'', pelo custo das falhas que ela causou.
      \end{historia}
      \begin{definicao}[\emph{Null safety}]
        O sistema de tipos separa tipos que \textbf{podem} ser nulos (\cod{String?}) dos que \textbf{não podem} (\cod{String}). Usar um valor possivelmente nulo sem tratar vira erro de compilação, não falha em produção.
      \end{definicao}
    \end{column}
    \begin{column}{0.48\textwidth}
      \centering
      \begin{tikzpicture}[node distance=3mm,every node/.style={font=\scriptsize,align=center}]
        \node[caixa,minimum width=2.2cm] (a) {\cod{String?}\\ \tiny pode ser nulo};
        \node[caixapetroleo,right=8mm of a,minimum width=2.2cm] (b) {\cod{String}\\ \tiny nunca nulo};
        \draw[seta] (a) -- node[rotulo,above]{\tiny verificação} (b);
        \node[below=3mm of a,font=\tiny,text width=2.2cm,align=center] {\cod{if (x != null)}\\ \cod{x ?: padrão}};
      \end{tikzpicture}
      \begin{pordentro}[Análise de fluxo]
        Depois de um \cod{if (x != null)}, o compilador sabe que \cod{x} não é nulo naquele trecho e permite usá-lo diretamente (\emph{smart cast} no Kotlin, promoção de tipo no Dart).
      \end{pordentro}
      \begin{atencao}
        O operador \cod{!!} (Kotlin) ou \cod{!} (Dart) desliga a proteção: só use quando tiver certeza.
      \end{atencao}
    \end{column}
  \end{columns}
\end{frame}

\begin{frame}{Os pilares da orientação a objetos}
  \begin{columns}[T]
    \begin{column}{0.54\textwidth}
      \scriptsize
      \begin{tabularx}{\linewidth}{@{}l L@{}}
        \toprule
        \cabecalho{Pilar} & \cabecalho{Significado e exemplo} \\
        \midrule
        Encapsulamento & esconder o estado interno e expor operações: a \cod{Reserva} só muda pelo método \cod{cancelar()} \\
        Abstração & mostrar o essencial e ocultar detalhes: \cod{FonteLeituras} sem dizer se é BLE ou MQTT \\
        Herança & uma classe especializa outra: \cod{SensorTemperatura} herda de \cod{Sensor} \\
        Polimorfismo & o mesmo comando, comportamentos diferentes: \cod{fonte.leituras()} em qualquer implementação \\
        \bottomrule
      \end{tabularx}
    \end{column}
    \begin{column}{0.44\textwidth}
      \begin{definicao}[Objeto e classe]
        A \textbf{classe} é o molde (atributos e métodos); o \textbf{objeto} é uma instância concreta, com estado próprio.
      \end{definicao}
      \begin{dica}[Composição em vez de herança]
        Prefira montar objetos com outros objetos (``tem um'') a criar hierarquias profundas (``é um''): mudanças ficam localizadas e o código, mais flexível.
      \end{dica}
    \end{column}
  \end{columns}
  \begin{analogia}
    Uma tomada é uma abstração: qualquer aparelho com o plugue certo funciona (polimorfismo), e ninguém precisa conhecer a fiação dentro da parede (encapsulamento).
  \end{analogia}
\end{frame}

\begin{frame}{Programação funcional}
  \begin{columns}[T]
    \begin{column}{0.5\textwidth}
      \begin{definicao}[Função pura]
        Para a mesma entrada, devolve sempre a mesma saída e não causa efeitos colaterais (não altera variáveis externas, não grava em disco, não acessa a rede).
      \end{definicao}
      \begin{definicao}[Funções de primeira classe e de ordem superior]
        Funções são valores: podem ser guardadas em variáveis, passadas como argumento e devolvidas. Uma função que recebe ou devolve outra é de \textbf{ordem superior} (\cod{map}, \cod{filter}).
      \end{definicao}
    \end{column}
    \begin{column}{0.48\textwidth}
      \centering
      \begin{tikzpicture}[node distance=4mm,every node/.style={font=\scriptsize,align=center}]
        \node[caixa,minimum width=1.9cm] (l) {leituras};
        \node[caixapetroleo,right=of l,minimum width=1.6cm] (f) {\cod{filter}\\ \tiny > 30\,°C};
        \node[caixapetroleo,right=of f,minimum width=1.6cm] (m) {\cod{map}\\ \tiny para alerta};
        \node[caixaambar,below=4mm of m,minimum width=1.6cm] (r) {alertas};
        \draw[seta] (l) -- (f); \draw[seta] (f) -- (m); \draw[seta] (m) -- (r);
      \end{tikzpicture}
      \begin{conceito}[Por que importa no mobile]
        Funções puras são fáceis de testar e seguras entre \emph{threads}; transformar listas com \cod{map}/\cod{filter} é mais legível que laços; interfaces declarativas tratam a tela como uma função do estado (Módulo 05).
      \end{conceito}
    \end{column}
  \end{columns}
\end{frame}

\begin{frame}{Memória: pilha, \emph{heap} e coleta de lixo}
  \begin{columns}[c]
    \begin{column}{0.46\textwidth}
      \centering
      \begin{tikzpicture}[every node/.style={font=\tiny}]
        \node[font=\scriptsize\bfseries,text=azulpetroleo] at (0.9,2.7) {Pilha};
        \foreach \y/\t in {2.2/main(),1.7/carregar(),1.2/calcular()} \node[caixa,minimum width=1.8cm] at (0.9,\y) {\t};
        \node[font=\scriptsize\bfseries,text=azulpetroleo] at (3.7,2.7) {\emph{Heap}};
        \draw[cinzamedio,rounded corners] (2.6,0.3) rectangle (4.8,2.45);
        \node[caixapetroleo,minimum width=1.3cm] (o1) at (3.3,1.95) {Bancada};
        \node[caixapetroleo,minimum width=1.3cm] (o2) at (4.1,1.2) {Lista};
        \node[caixa,minimum width=1.3cm,fill=cinzaclaro] (o3) at (3.3,0.65) {sem ref.};
        \draw[seta] (1.8,2.2) -- (o1); \draw[seta] (1.8,1.7) -- (o2); \draw[seta] (o2) -- (o1);
        \node[font=\tiny,text=vermelhoerro] at (3.7,0.2) {coletado};
      \end{tikzpicture}
    \end{column}
    \begin{column}{0.52\textwidth}
      \begin{definicao}[Pilha e \emph{heap}]
        A \textbf{pilha} guarda as chamadas de função e suas variáveis locais, liberadas ao retornar. O \textbf{\emph{heap}} guarda os objetos, que vivem enquanto houver referências a eles.
      \end{definicao}
      \begin{definicao}[Coleta de lixo]
        O \emph{runtime} encontra objetos que não são mais alcançáveis a partir das referências vivas e libera sua memória automaticamente.
      \end{definicao}
      \begin{pordentro}[Gerações]
        A maioria dos objetos morre jovem. Por isso o ART e a VM do Dart separam objetos novos e antigos e coletam os novos com frequência e rapidez.
      \end{pordentro}
    \end{column}
  \end{columns}
\end{frame}

\begin{frame}{Imutabilidade e igualdade}
  \begin{columns}[T]
    \begin{column}{0.5\textwidth}
      \begin{definicao}[Imutável]
        Objeto que não muda depois de criado. Para ``alterar'', cria-se uma cópia com o campo novo (\cod{copy} no Kotlin, \cod{copyWith} no Dart).
      \end{definicao}
      \begin{conceito}[Vantagens]
        \begin{itemize}
          \item Ninguém altera o objeto ``pelas costas''.
          \item Seguro para compartilhar entre \emph{threads} e \emph{isolates}.
          \item Comparar o antes e o depois de um estado fica trivial (Módulo 07).
        \end{itemize}
      \end{conceito}
    \end{column}
    \begin{column}{0.48\textwidth}
      \begin{definicao}[Igualdade de valor e de identidade]
        \textbf{Identidade}: as duas variáveis apontam para o \textbf{mesmo} objeto. \textbf{Valor}: objetos diferentes com o mesmo conteúdo. \cod{data class} (Kotlin) e \emph{records} (Dart) comparam por valor.
      \end{definicao}
      \begin{erro}[Lista mutável compartilhada]
        Dois componentes recebem a mesma lista; um a altera e o outro mostra dados que não esperava. Exponha listas somente leitura.
      \end{erro}
    \end{column}
  \end{columns}
\end{frame}

% =============================================================================
\section{Fundamentos: linguagens de programação}

\begin{frame}{Paradigmas de programação}
  \begin{columns}[T]
    \begin{column}{0.54\textwidth}
      \scriptsize
      \begin{tabularx}{\linewidth}{@{}l L@{}}
        \toprule
        \cabecalho{Paradigma} & \cabecalho{Ideia central} \\
        \midrule
        Imperativo & sequência de comandos que alteram o estado (\emph{faça isto, depois aquilo}) \\
        Orientado a objetos & dados e comportamentos agrupados em objetos que trocam mensagens \\
        Funcional & funções puras, valores imutáveis, composição de funções \\
        Declarativo & descreve \textbf{o que} se quer, não \textbf{como} obter (SQL, interfaces em Compose e Flutter) \\
        \bottomrule
      \end{tabularx}
    \end{column}
    \begin{column}{0.44\textwidth}
      \begin{definicao}[Paradigma]
        Estilo de organizar o raciocínio e o código. Linguagens modernas são \textbf{multiparadigma}: Kotlin e Dart misturam imperativo, objetos e funcional.
      \end{definicao}
      \begin{dica}
        Use o paradigma que deixa cada parte mais clara: objetos para modelar o domínio, funções para transformar dados, declarativo para descrever a interface.
      \end{dica}
    \end{column}
  \end{columns}
  \begin{historia}
    Fortran (1957) era imperativo; Lisp (1958), funcional; Simula (1967) e Smalltalk (anos 1970) criaram a orientação a objetos. Kotlin (2011, JetBrains) e Dart (2011, Google) herdaram ideias de todos.
  \end{historia}
\end{frame}

\begin{frame}{Sistemas de tipos}
  \begin{columns}[T]
    \begin{column}{0.5\textwidth}
      \begin{definicao}[Tipo]
        Conjunto de valores possíveis e das operações permitidas sobre eles. Somar dois \cod{Int} faz sentido; somar um \cod{Int} a uma \cod{Bancada}, não.
      \end{definicao}
      \scriptsize
      \begin{tabularx}{\linewidth}{@{}l L@{}}
        \toprule
        \cabecalho{Classificação} & \cabecalho{Significado} \\
        \midrule
        Estática & tipos verificados na compilação (Kotlin, Dart, Java) \\
        Dinâmica & tipos verificados na execução (Python, JavaScript) \\
        Forte & conversões perigosas são proibidas ou explícitas \\
        Inferência & o compilador deduz o tipo (\cod{val x = 3}) \\
        \bottomrule
      \end{tabularx}
    \end{column}
    \begin{column}{0.48\textwidth}
      \begin{conceito}[Por que tipagem estática ajuda]
        Muitos erros aparecem \textbf{antes} de o app rodar; a IDE completa código e refatora com segurança; o tipo documenta a intenção.
      \end{conceito}
      \begin{pordentro}[Solidez (\emph{soundness})]
        Um sistema de tipos é sólido quando garante que, se o programa compila, uma variável nunca terá em execução um valor de outro tipo. O Dart 3 é sólido, inclusive quanto a nulos.
      \end{pordentro}
    \end{column}
  \end{columns}
\end{frame}

\begin{frame}{O ``erro de um bilhão de dólares'': o nulo}
  \begin{columns}[T]
    \begin{column}{0.5\textwidth}
      \begin{historia}
        Tony Hoare introduziu a referência nula na linguagem ALGOL W em 1965. Décadas depois, chamou-a de seu ``erro de um bilhão de dólares'', pelo custo das falhas que ela causou.
      \end{historia}
      \begin{definicao}[\emph{Null safety}]
        O sistema de tipos separa tipos que \textbf{podem} ser nulos (\cod{String?}) dos que \textbf{não podem} (\cod{String}). Usar um valor possivelmente nulo sem tratar vira erro de compilação, não falha em produção.
      \end{definicao}
    \end{column}
    \begin{column}{0.48\textwidth}
      \centering
      \begin{tikzpicture}[node distance=3mm,every node/.style={font=\scriptsize,align=center}]
        \node[caixa,minimum width=2.2cm] (a) {\cod{String?}\\ \tiny pode ser nulo};
        \node[caixapetroleo,right=8mm of a,minimum width=2.2cm] (b) {\cod{String}\\ \tiny nunca nulo};
        \draw[seta] (a) -- node[rotulo,above]{\tiny verificação} (b);
        \node[below=3mm of a,font=\tiny,text width=2.2cm,align=center] {\cod{if (x != null)}\\ \cod{x ?: padrão}};
      \end{tikzpicture}
      \begin{pordentro}[Análise de fluxo]
        Depois de um \cod{if (x != null)}, o compilador sabe que \cod{x} não é nulo naquele trecho e permite usá-lo diretamente (\emph{smart cast} no Kotlin, promoção de tipo no Dart).
      \end{pordentro}
      \begin{atencao}
        O operador \cod{!!} (Kotlin) ou \cod{!} (Dart) desliga a proteção: só use quando tiver certeza.
      \end{atencao}
    \end{column}
  \end{columns}
\end{frame}

\begin{frame}{Os pilares da orientação a objetos}
  \begin{columns}[T]
    \begin{column}{0.54\textwidth}
      \scriptsize
      \begin{tabularx}{\linewidth}{@{}l L@{}}
        \toprule
        \cabecalho{Pilar} & \cabecalho{Significado e exemplo} \\
        \midrule
        Encapsulamento & esconder o estado interno e expor operações: a \cod{Reserva} só muda pelo método \cod{cancelar()} \\
        Abstração & mostrar o essencial e ocultar detalhes: \cod{FonteLeituras} sem dizer se é BLE ou MQTT \\
        Herança & uma classe especializa outra: \cod{SensorTemperatura} herda de \cod{Sensor} \\
        Polimorfismo & o mesmo comando, comportamentos diferentes: \cod{fonte.leituras()} em qualquer implementação \\
        \bottomrule
      \end{tabularx}
    \end{column}
    \begin{column}{0.44\textwidth}
      \begin{definicao}[Objeto e classe]
        A \textbf{classe} é o molde (atributos e métodos); o \textbf{objeto} é uma instância concreta, com estado próprio.
      \end{definicao}
      \begin{dica}[Composição em vez de herança]
        Prefira montar objetos com outros objetos (``tem um'') a criar hierarquias profundas (``é um''): mudanças ficam localizadas e o código, mais flexível.
      \end{dica}
    \end{column}
  \end{columns}
  \begin{analogia}
    Uma tomada é uma abstração: qualquer aparelho com o plugue certo funciona (polimorfismo), e ninguém precisa conhecer a fiação dentro da parede (encapsulamento).
  \end{analogia}
\end{frame}

\begin{frame}{Programação funcional}
  \begin{columns}[T]
    \begin{column}{0.5\textwidth}
      \begin{definicao}[Função pura]
        Para a mesma entrada, devolve sempre a mesma saída e não causa efeitos colaterais (não altera variáveis externas, não grava em disco, não acessa a rede).
      \end{definicao}
      \begin{definicao}[Funções de primeira classe e de ordem superior]
        Funções são valores: podem ser guardadas em variáveis, passadas como argumento e devolvidas. Uma função que recebe ou devolve outra é de \textbf{ordem superior} (\cod{map}, \cod{filter}).
      \end{definicao}
    \end{column}
    \begin{column}{0.48\textwidth}
      \centering
      \begin{tikzpicture}[node distance=4mm,every node/.style={font=\scriptsize,align=center}]
        \node[caixa,minimum width=1.9cm] (l) {leituras};
        \node[caixapetroleo,right=of l,minimum width=1.6cm] (f) {\cod{filter}\\ \tiny > 30\,°C};
        \node[caixapetroleo,right=of f,minimum width=1.6cm] (m) {\cod{map}\\ \tiny para alerta};
        \node[caixaambar,below=4mm of m,minimum width=1.6cm] (r) {alertas};
        \draw[seta] (l) -- (f); \draw[seta] (f) -- (m); \draw[seta] (m) -- (r);
      \end{tikzpicture}
      \begin{conceito}[Por que importa no mobile]
        Funções puras são fáceis de testar e seguras entre \emph{threads}; transformar listas com \cod{map}/\cod{filter} é mais legível que laços; interfaces declarativas tratam a tela como uma função do estado (Módulo 05).
      \end{conceito}
    \end{column}
  \end{columns}
\end{frame}

\begin{frame}{Memória: pilha, \emph{heap} e coleta de lixo}
  \begin{columns}[c]
    \begin{column}{0.46\textwidth}
      \centering
      \begin{tikzpicture}[every node/.style={font=\tiny}]
        \node[font=\scriptsize\bfseries,text=azulpetroleo] at (0.9,2.7) {Pilha};
        \foreach \y/\t in {2.2/main(),1.7/carregar(),1.2/calcular()} \node[caixa,minimum width=1.8cm] at (0.9,\y) {\t};
        \node[font=\scriptsize\bfseries,text=azulpetroleo] at (3.7,2.7) {\emph{Heap}};
        \draw[cinzamedio,rounded corners] (2.6,0.3) rectangle (4.8,2.45);
        \node[caixapetroleo,minimum width=1.3cm] (o1) at (3.3,1.95) {Bancada};
        \node[caixapetroleo,minimum width=1.3cm] (o2) at (4.1,1.2) {Lista};
        \node[caixa,minimum width=1.3cm,fill=cinzaclaro] (o3) at (3.3,0.65) {sem ref.};
        \draw[seta] (1.8,2.2) -- (o1); \draw[seta] (1.8,1.7) -- (o2); \draw[seta] (o2) -- (o1);
        \node[font=\tiny,text=vermelhoerro] at (3.7,0.2) {coletado};
      \end{tikzpicture}
    \end{column}
    \begin{column}{0.52\textwidth}
      \begin{definicao}[Pilha e \emph{heap}]
        A \textbf{pilha} guarda as chamadas de função e suas variáveis locais, liberadas ao retornar. O \textbf{\emph{heap}} guarda os objetos, que vivem enquanto houver referências a eles.
      \end{definicao}
      \begin{definicao}[Coleta de lixo]
        O \emph{runtime} encontra objetos que não são mais alcançáveis a partir das referências vivas e libera sua memória automaticamente.
      \end{definicao}
      \begin{pordentro}[Gerações]
        A maioria dos objetos morre jovem. Por isso o ART e a VM do Dart separam objetos novos e antigos e coletam os novos com frequência e rapidez.
      \end{pordentro}
    \end{column}
  \end{columns}
\end{frame}

\begin{frame}{Imutabilidade e igualdade}
  \begin{columns}[T]
    \begin{column}{0.5\textwidth}
      \begin{definicao}[Imutável]
        Objeto que não muda depois de criado. Para ``alterar'', cria-se uma cópia com o campo novo (\cod{copy} no Kotlin, \cod{copyWith} no Dart).
      \end{definicao}
      \begin{conceito}[Vantagens]
        \begin{itemize}
          \item Ninguém altera o objeto ``pelas costas''.
          \item Seguro para compartilhar entre \emph{threads} e \emph{isolates}.
          \item Comparar o antes e o depois de um estado fica trivial (Módulo 07).
        \end{itemize}
      \end{conceito}
    \end{column}
    \begin{column}{0.48\textwidth}
      \begin{definicao}[Igualdade de valor e de identidade]
        \textbf{Identidade}: as duas variáveis apontam para o \textbf{mesmo} objeto. \textbf{Valor}: objetos diferentes com o mesmo conteúdo. \cod{data class} (Kotlin) e \emph{records} (Dart) comparam por valor.
      \end{definicao}
      \begin{erro}[Lista mutável compartilhada]
        Dois componentes recebem a mesma lista; um a altera e o outro mostra dados que não esperava. Exponha listas somente leitura.
      \end{erro}
    \end{column}
  \end{columns}
\end{frame}

% =============================================================================
\section{Fundamentos: linguagens de programação}

\begin{frame}{Paradigmas de programação}
  \begin{columns}[T]
    \begin{column}{0.54\textwidth}
      \scriptsize
      \begin{tabularx}{\linewidth}{@{}l L@{}}
        \toprule
        \cabecalho{Paradigma} & \cabecalho{Ideia central} \\
        \midrule
        Imperativo & sequência de comandos que alteram o estado (\emph{faça isto, depois aquilo}) \\
        Orientado a objetos & dados e comportamentos agrupados em objetos que trocam mensagens \\
        Funcional & funções puras, valores imutáveis, composição de funções \\
        Declarativo & descreve \textbf{o que} se quer, não \textbf{como} obter (SQL, interfaces em Compose e Flutter) \\
        \bottomrule
      \end{tabularx}
    \end{column}
    \begin{column}{0.44\textwidth}
      \begin{definicao}[Paradigma]
        Estilo de organizar o raciocínio e o código. Linguagens modernas são \textbf{multiparadigma}: Kotlin e Dart misturam imperativo, objetos e funcional.
      \end{definicao}
      \begin{dica}
        Use o paradigma que deixa cada parte mais clara: objetos para modelar o domínio, funções para transformar dados, declarativo para descrever a interface.
      \end{dica}
    \end{column}
  \end{columns}
  \begin{historia}
    Fortran (1957) era imperativo; Lisp (1958), funcional; Simula (1967) e Smalltalk (anos 1970) criaram a orientação a objetos. Kotlin (2011, JetBrains) e Dart (2011, Google) herdaram ideias de todos.
  \end{historia}
\end{frame}

\begin{frame}{Sistemas de tipos}
  \begin{columns}[T]
    \begin{column}{0.5\textwidth}
      \begin{definicao}[Tipo]
        Conjunto de valores possíveis e das operações permitidas sobre eles. Somar dois \cod{Int} faz sentido; somar um \cod{Int} a uma \cod{Bancada}, não.
      \end{definicao}
      \scriptsize
      \begin{tabularx}{\linewidth}{@{}l L@{}}
        \toprule
        \cabecalho{Classificação} & \cabecalho{Significado} \\
        \midrule
        Estática & tipos verificados na compilação (Kotlin, Dart, Java) \\
        Dinâmica & tipos verificados na execução (Python, JavaScript) \\
        Forte & conversões perigosas são proibidas ou explícitas \\
        Inferência & o compilador deduz o tipo (\cod{val x = 3}) \\
        \bottomrule
      \end{tabularx}
    \end{column}
    \begin{column}{0.48\textwidth}
      \begin{conceito}[Por que tipagem estática ajuda]
        Muitos erros aparecem \textbf{antes} de o app rodar; a IDE completa código e refatora com segurança; o tipo documenta a intenção.
      \end{conceito}
      \begin{pordentro}[Solidez (\emph{soundness})]
        Um sistema de tipos é sólido quando garante que, se o programa compila, uma variável nunca terá em execução um valor de outro tipo. O Dart 3 é sólido, inclusive quanto a nulos.
      \end{pordentro}
    \end{column}
  \end{columns}
\end{frame}

\begin{frame}{O ``erro de um bilhão de dólares'': o nulo}
  \begin{columns}[T]
    \begin{column}{0.5\textwidth}
      \begin{historia}
        Tony Hoare introduziu a referência nula na linguagem ALGOL W em 1965. Décadas depois, chamou-a de seu ``erro de um bilhão de dólares'', pelo custo das falhas que ela causou.
      \end{historia}
      \begin{definicao}[\emph{Null safety}]
        O sistema de tipos separa tipos que \textbf{podem} ser nulos (\cod{String?}) dos que \textbf{não podem} (\cod{String}). Usar um valor possivelmente nulo sem tratar vira erro de compilação, não falha em produção.
      \end{definicao}
    \end{column}
    \begin{column}{0.48\textwidth}
      \centering
      \begin{tikzpicture}[node distance=3mm,every node/.style={font=\scriptsize,align=center}]
        \node[caixa,minimum width=2.2cm] (a) {\cod{String?}\\ \tiny pode ser nulo};
        \node[caixapetroleo,right=8mm of a,minimum width=2.2cm] (b) {\cod{String}\\ \tiny nunca nulo};
        \draw[seta] (a) -- node[rotulo,above]{\tiny verificação} (b);
        \node[below=3mm of a,font=\tiny,text width=2.2cm,align=center] {\cod{if (x != null)}\\ \cod{x ?: padrão}};
      \end{tikzpicture}
      \begin{pordentro}[Análise de fluxo]
        Depois de um \cod{if (x != null)}, o compilador sabe que \cod{x} não é nulo naquele trecho e permite usá-lo diretamente (\emph{smart cast} no Kotlin, promoção de tipo no Dart).
      \end{pordentro}
      \begin{atencao}
        O operador \cod{!!} (Kotlin) ou \cod{!} (Dart) desliga a proteção: só use quando tiver certeza.
      \end{atencao}
    \end{column}
  \end{columns}
\end{frame}

\begin{frame}{Os pilares da orientação a objetos}
  \begin{columns}[T]
    \begin{column}{0.54\textwidth}
      \scriptsize
      \begin{tabularx}{\linewidth}{@{}l L@{}}
        \toprule
        \cabecalho{Pilar} & \cabecalho{Significado e exemplo} \\
        \midrule
        Encapsulamento & esconder o estado interno e expor operações: a \cod{Reserva} só muda pelo método \cod{cancelar()} \\
        Abstração & mostrar o essencial e ocultar detalhes: \cod{FonteLeituras} sem dizer se é BLE ou MQTT \\
        Herança & uma classe especializa outra: \cod{SensorTemperatura} herda de \cod{Sensor} \\
        Polimorfismo & o mesmo comando, comportamentos diferentes: \cod{fonte.leituras()} em qualquer implementação \\
        \bottomrule
      \end{tabularx}
    \end{column}
    \begin{column}{0.44\textwidth}
      \begin{definicao}[Objeto e classe]
        A \textbf{classe} é o molde (atributos e métodos); o \textbf{objeto} é uma instância concreta, com estado próprio.
      \end{definicao}
      \begin{dica}[Composição em vez de herança]
        Prefira montar objetos com outros objetos (``tem um'') a criar hierarquias profundas (``é um''): mudanças ficam localizadas e o código, mais flexível.
      \end{dica}
    \end{column}
  \end{columns}
  \begin{analogia}
    Uma tomada é uma abstração: qualquer aparelho com o plugue certo funciona (polimorfismo), e ninguém precisa conhecer a fiação dentro da parede (encapsulamento).
  \end{analogia}
\end{frame}

\begin{frame}{Programação funcional}
  \begin{columns}[T]
    \begin{column}{0.5\textwidth}
      \begin{definicao}[Função pura]
        Para a mesma entrada, devolve sempre a mesma saída e não causa efeitos colaterais (não altera variáveis externas, não grava em disco, não acessa a rede).
      \end{definicao}
      \begin{definicao}[Funções de primeira classe e de ordem superior]
        Funções são valores: podem ser guardadas em variáveis, passadas como argumento e devolvidas. Uma função que recebe ou devolve outra é de \textbf{ordem superior} (\cod{map}, \cod{filter}).
      \end{definicao}
    \end{column}
    \begin{column}{0.48\textwidth}
      \centering
      \begin{tikzpicture}[node distance=4mm,every node/.style={font=\scriptsize,align=center}]
        \node[caixa,minimum width=1.9cm] (l) {leituras};
        \node[caixapetroleo,right=of l,minimum width=1.6cm] (f) {\cod{filter}\\ \tiny > 30\,°C};
        \node[caixapetroleo,right=of f,minimum width=1.6cm] (m) {\cod{map}\\ \tiny para alerta};
        \node[caixaambar,below=4mm of m,minimum width=1.6cm] (r) {alertas};
        \draw[seta] (l) -- (f); \draw[seta] (f) -- (m); \draw[seta] (m) -- (r);
      \end{tikzpicture}
      \begin{conceito}[Por que importa no mobile]
        Funções puras são fáceis de testar e seguras entre \emph{threads}; transformar listas com \cod{map}/\cod{filter} é mais legível que laços; interfaces declarativas tratam a tela como uma função do estado (Módulo 05).
      \end{conceito}
    \end{column}
  \end{columns}
\end{frame}

\begin{frame}{Memória: pilha, \emph{heap} e coleta de lixo}
  \begin{columns}[c]
    \begin{column}{0.46\textwidth}
      \centering
      \begin{tikzpicture}[every node/.style={font=\tiny}]
        \node[font=\scriptsize\bfseries,text=azulpetroleo] at (0.9,2.7) {Pilha};
        \foreach \y/\t in {2.2/main(),1.7/carregar(),1.2/calcular()} \node[caixa,minimum width=1.8cm] at (0.9,\y) {\t};
        \node[font=\scriptsize\bfseries,text=azulpetroleo] at (3.7,2.7) {\emph{Heap}};
        \draw[cinzamedio,rounded corners] (2.6,0.3) rectangle (4.8,2.45);
        \node[caixapetroleo,minimum width=1.3cm] (o1) at (3.3,1.95) {Bancada};
        \node[caixapetroleo,minimum width=1.3cm] (o2) at (4.1,1.2) {Lista};
        \node[caixa,minimum width=1.3cm,fill=cinzaclaro] (o3) at (3.3,0.65) {sem ref.};
        \draw[seta] (1.8,2.2) -- (o1); \draw[seta] (1.8,1.7) -- (o2); \draw[seta] (o2) -- (o1);
        \node[font=\tiny,text=vermelhoerro] at (3.7,0.2) {coletado};
      \end{tikzpicture}
    \end{column}
    \begin{column}{0.52\textwidth}
      \begin{definicao}[Pilha e \emph{heap}]
        A \textbf{pilha} guarda as chamadas de função e suas variáveis locais, liberadas ao retornar. O \textbf{\emph{heap}} guarda os objetos, que vivem enquanto houver referências a eles.
      \end{definicao}
      \begin{definicao}[Coleta de lixo]
        O \emph{runtime} encontra objetos que não são mais alcançáveis a partir das referências vivas e libera sua memória automaticamente.
      \end{definicao}
      \begin{pordentro}[Gerações]
        A maioria dos objetos morre jovem. Por isso o ART e a VM do Dart separam objetos novos e antigos e coletam os novos com frequência e rapidez.
      \end{pordentro}
    \end{column}
  \end{columns}
\end{frame}

\begin{frame}{Imutabilidade e igualdade}
  \begin{columns}[T]
    \begin{column}{0.5\textwidth}
      \begin{definicao}[Imutável]
        Objeto que não muda depois de criado. Para ``alterar'', cria-se uma cópia com o campo novo (\cod{copy} no Kotlin, \cod{copyWith} no Dart).
      \end{definicao}
      \begin{conceito}[Vantagens]
        \begin{itemize}
          \item Ninguém altera o objeto ``pelas costas''.
          \item Seguro para compartilhar entre \emph{threads} e \emph{isolates}.
          \item Comparar o antes e o depois de um estado fica trivial (Módulo 07).
        \end{itemize}
      \end{conceito}
    \end{column}
    \begin{column}{0.48\textwidth}
      \begin{definicao}[Igualdade de valor e de identidade]
        \textbf{Identidade}: as duas variáveis apontam para o \textbf{mesmo} objeto. \textbf{Valor}: objetos diferentes com o mesmo conteúdo. \cod{data class} (Kotlin) e \emph{records} (Dart) comparam por valor.
      \end{definicao}
      \begin{erro}[Lista mutável compartilhada]
        Dois componentes recebem a mesma lista; um a altera e o outro mostra dados que não esperava. Exponha listas somente leitura.
      \end{erro}
    \end{column}
  \end{columns}
\end{frame}


% =============================================================================
\section{Primeiros passos}

\begin{frame}{Por que Kotlin?}
  \begin{columns}[T]
    \begin{column}{0.55\textwidth}
      \begin{itemize}
        \item Criada pela JetBrains; linguagem preferencial do Android desde 2019.
        \item \textbf{Concisa}: menos código repetitivo que Java.
        \item \textbf{Segura}: o sistema de tipos distingue o que pode e o que não pode ser nulo.
        \item \textbf{Interoperável}: usa qualquer biblioteca Java.
        \item \textbf{Multiplataforma}: JVM, Android, iOS, web e desktop (Módulo 19).
        \item Suporte nativo a programação assíncrona com \emph{coroutines} (Módulo 06).
      \end{itemize}
    \end{column}
    \begin{column}{0.43\textwidth}
      \begin{exemplo}[Java]
        \ttfamily\scriptsize
        public class Sensor \{\\
        \ \ private final String nome;\\
        \ \ public Sensor(String nome) \{\\
        \ \ \ \ this.nome = nome;\\
        \ \ \}\\
        \ \ public String getNome() \{\\
        \ \ \ \ return nome;\\
        \ \ \}\\
        \ \ // + equals, hashCode, toString\ldots\\
        \}
      \end{exemplo}
      \begin{exemplo}[Kotlin]
        \ttfamily\scriptsize data class Sensor(val nome: String)
      \end{exemplo}
    \end{column}
  \end{columns}
\end{frame}

\begin{frame}[fragile]{Variáveis e tipos básicos}
  \begin{columns}[T]
    \begin{column}{0.54\textwidth}
\begin{codigo}[title=Variaveis.kt]{Kotlin}
fun main() {
  val bancada = 3          // imutável (prefira!)
  var temperatura = 24.5   // mutável
  temperatura += 0.3

  val nome: String = "Osciloscópio"  // tipo explícito
  val ligado = true                  // Boolean
  val leituras: Long = 1_000_000L    // separador _
  val letra: Char = 'A'

  // String templates
  println("Bancada $bancada: $temperatura °C")
  println("Nome tem ${nome.length} letras")

  // Texto multilinha
  val relatorio = """
      |Equipamento: $nome
      |Estado: ${if (ligado) "ON" else "OFF"}
  """.trimMargin()
}
\end{codigo}
    \end{column}
    \begin{column}{0.44\textwidth}
      \begin{conceito}[\cod{val} \emph{versus} \cod{var}]
        \cod{val} cria uma referência que não pode ser reatribuída. \cod{var} pode. Use \cod{val} por padrão e só troque se precisar.
      \end{conceito}
      \footnotesize
      \begin{tabular}{@{}l l@{}}
        \toprule
        \cabecalho{Tipo} & \cabecalho{Exemplo} \\
        \midrule
        \cod{Int} / \cod{Long} & \cod{42} / \cod{42L} \\
        \cod{Double} / \cod{Float} & \cod{3.14} / \cod{3.14f} \\
        \cod{Boolean} & \cod{true} / \cod{false} \\
        \cod{Char} / \cod{String} & \cod{'a'} / \cod{"texto"} \\
        \bottomrule
      \end{tabular}
      \begin{dica}
        Conversões são explícitas: \cod{x.toInt()}, \cod{"12".toDouble()}, \cod{"abc".toIntOrNull()}.
      \end{dica}
    \end{column}
  \end{columns}
\end{frame}

% =============================================================================
\section{Controle de fluxo}

\begin{frame}[fragile]{\cod{if} e \cod{when} são expressões}
  \begin{columns}[T]
    \begin{column}{0.49\textwidth}
\begin{codigo}[title=if como expressão]{Kotlin}
val status = if (temp > 80.0) "ALERTA" else "OK"

val nivel = if (temp < 20) {
  "frio"
} else if (temp < 30) {
  "normal"
} else {
  "quente"
}
\end{codigo}
      \begin{dica}
        Kotlin não tem operador ternário (\cod{a ? b : c}): o \cod{if} já devolve valor.
      \end{dica}
    \end{column}
    \begin{column}{0.49\textwidth}
\begin{codigo}[title=when: o switch turbinado]{Kotlin}
fun classificar(t: Double) = when {
  t < 0      -> "abaixo da faixa"
  t in 0.0..35.0 -> "normal"
  else       -> "acima da faixa"
}

fun descrever(x: Any): String = when (x) {
  is Int    -> "inteiro $x"
  is String -> "texto de ${x.length} letras"
  1, 2, 3   -> "pequeno"
  else      -> "desconhecido"
}
\end{codigo}
    \end{column}
  \end{columns}
\end{frame}

\begin{frame}[fragile]{Ranges e laços}
  \begin{columns}[T]
    \begin{column}{0.52\textwidth}
\begin{codigo}[title=Laços]{Kotlin}
for (i in 1..5) print(i)          // 1 2 3 4 5
for (i in 0 until 5) print(i)     // 0 1 2 3 4
for (i in 10 downTo 0 step 2) print(i)
for (i in 0..<5) print(i)         // igual a until

val sensores = listOf("DHT22", "BMP280", "LDR")
for (s in sensores) println(s)
for ((i, s) in sensores.withIndex()) println("$i: $s")

var tentativas = 0
while (tentativas < 3) {
  tentativas++
}
repeat(3) { println("ping $it") }
\end{codigo}
    \end{column}
    \begin{column}{0.46\textwidth}
      \begin{conceito}[Ranges]
        \cod{a..b} inclui \cod{b}; \cod{a until b} e \cod{a..<b} excluem \cod{b}. O operador \cod{in} testa se um valor pertence ao intervalo:\\[2pt]
        \cod{if (tensao in 4.75..5.25) \{ \ldots \}}
      \end{conceito}
      \begin{dica}
        Em código Android você raramente escreverá laços: operações como \cod{map} e \cod{filter} (seção 6) costumam ser mais claras.
      \end{dica}
    \end{column}
  \end{columns}
\end{frame}

% =============================================================================
\section{Null safety}

\begin{frame}[fragile]{Tipos anuláveis}
  \begin{columns}[T]
    \begin{column}{0.52\textwidth}
\begin{codigo}[title=Null safety]{Kotlin}
var nome: String = "LDR"
// nome = null          // ERRO de compilação

var apelido: String? = null   // pode ser nulo
// apelido.length       // ERRO: pode ser nulo
// 1) chamada segura: devolve null se for null
val tam: Int? = apelido?.length
// 2) elvis: valor padrão quando null
val tam2: Int = apelido?.length ?: 0

// 3) smart cast após checagem
if (apelido != null) {
  println(apelido.length)    // aqui é String
}

// 4) let: executa só se não for null
apelido?.let { println("Apelido: $it") }

// 5) !! afirma que não é null (evite!)
val tam3 = apelido!!.length  // lança NPE se null
\end{codigo}
    \end{column}
    \begin{column}{0.46\textwidth}
      \begin{conceito}[A ideia]
        O compilador sabe se um valor pode ser nulo. \cod{String} nunca é nulo; \cod{String?} pode ser. Você é \textbf{obrigado} a tratar o caso nulo antes de usar.
      \end{conceito}
      \begin{center}
      \begin{tikzpicture}[node distance=3mm]
        \node[caixaambar] (a) {\cod{apelido?.length ?: 0}};
        \node[caixa,below left=5mm and -8mm of a] (n) {null $\rightarrow$ \cod{0}};
        \node[caixa,below right=5mm and -8mm of a] (v) {"Lu" $\rightarrow$ \cod{2}};
        \draw[seta] (a) -- (n);
        \draw[seta] (a) -- (v);
      \end{tikzpicture}
      \end{center}
      \begin{erro}
        Usar \cod{!!} para ``calar'' o compilador traz de volta o \cod{NullPointerException}.
      \end{erro}
    \end{column}
  \end{columns}
\end{frame}

% =============================================================================
\section{Funções e lambdas}

\begin{frame}[fragile]{Funções}
  \begin{columns}[T]
    \begin{column}{0.53\textwidth}
\begin{codigo}[title=Funcoes.kt]{Kotlin}
// Forma completa
fun media(valores: List<Double>): Double {
  return valores.sum() / valores.size
}

// Corpo de expressão única
fun celsiusParaF(c: Double) = c * 9 / 5 + 32

// Parâmetros com valor padrão
fun conectar(
  host: String,
  porta: Int = 1883,
  tls: Boolean = false
) { /* ... */ }

conectar("broker.local")
conectar("broker.local", tls = true)  // nomeado

// Função sem retorno útil devolve Unit
fun log(msg: String) { println(msg) }
\end{codigo}
    \end{column}
    \begin{column}{0.45\textwidth}
\begin{codigo}[title=Funções de extensão]{Kotlin}
// "adiciona" um método a uma classe existente
fun Double.emCelsius() = "%.1f °C".format(this)

val MAC = Regex("([0-9A-F]{2}:){5}[0-9A-F]{2}")
fun String.ehMacValido() = matches(MAC)

println(24.567.emCelsius())   // 24.6 °C
println("24:6F:28:AB:CD:EF".ehMacValido())
\end{codigo}
      \begin{conceito}[Extensões]
        Não modificam a classe original: são funções comuns com uma sintaxe mais legível. São muito usadas no Android (\cod{Modifier.padding}, \cod{Context.getString}\ldots).
      \end{conceito}
    \end{column}
  \end{columns}
\end{frame}

\begin{frame}[fragile]{Lambdas e funções de ordem superior}
  \begin{columns}[T]
    \begin{column}{0.53\textwidth}
\begin{codigo}[title=Lambdas]{Kotlin}
// Lambda: função anônima atribuída a uma variável
val dobro: (Int) -> Int = { x -> x * 2 }
val soma = { a: Int, b: Int -> a + b }

// Com um único parâmetro, use "it"
val ehPar: (Int) -> Boolean = { it % 2 == 0 }

// Função de ordem superior: recebe uma função
fun medir(n: Int, acao: () -> Unit): Long {
  val inicio = System.nanoTime()
  repeat(n) { acao() }
  return System.nanoTime() - inicio
}

// "Trailing lambda": a última lambda sai dos ()
val tempo = medir(1000) {
  calcularFFT()
}
\end{codigo}
    \end{column}
    \begin{column}{0.45\textwidth}
      \begin{conceito}[Tipos de função]
        \cod{(Int) -> Int} lê-se ``recebe um \cod{Int} e devolve um \cod{Int}''. \cod{() -> Unit} não recebe nada e não devolve nada.
      \end{conceito}
      \begin{atencao}[Isto é Compose!]
        Toda a interface do Compose usa \emph{trailing lambdas}:\\[2pt]
        \cod{Button(onClick = \{ ligar() \}) \{}\\
        \cod{\ \ Text("Ligar")}\\
        \cod{\}}\\[2pt]
        \cod{onClick} recebe uma lambda e o conteúdo do botão é outra lambda.
      \end{atencao}
    \end{column}
  \end{columns}
\end{frame}

% =============================================================================
\section{Orientação a objetos}

\begin{frame}[fragile]{Classes, construtores e propriedades}
  \begin{columns}[T]
    \begin{column}{0.54\textwidth}
\begin{codigo}[title=Classes]{Kotlin}
class Equipamento(
  val patrimonio: String,       // propriedade imutável
  var descricao: String,        // propriedade mutável
  private var horasDeUso: Int = 0
) {
  val precisaManutencao: Boolean   // calculada
    get() = horasDeUso > 500

  init {
    require(patrimonio.isNotBlank()) { "Patrimônio obrigatório" }
  }

  fun registrarUso(horas: Int) { horasDeUso += horas }
}

val osc = Equipamento("IFMA-0042", "Osciloscópio")
osc.registrarUso(12)
println(osc.precisaManutencao)
\end{codigo}
    \end{column}
    \begin{column}{0.44\textwidth}
      \small
      \begin{itemize}
        \item O \termo{construtor primário} fica no cabeçalho da classe.
        \item \cod{val}/\cod{var} no construtor já criam as propriedades.
        \item \cod{init} roda na criação do objeto.
        \item Não existe \cod{new}.
        \item Visibilidade: \cod{public} (padrão), \cod{private}, \cod{protected}, \cod{internal} (mesmo módulo).
      \end{itemize}
      \begin{dica}
        \cod{require} valida argumentos (lança \cod{IllegalArgumentException}); \cod{check} valida estado (lança \cod{IllegalStateException}).
      \end{dica}
    \end{column}
  \end{columns}
\end{frame}

\begin{frame}[fragile]{Herança e interfaces}
  \begin{columns}[T]
    \begin{column}{0.53\textwidth}
\begin{codigo}[title=Herança]{Kotlin}
interface Medidor {
  val unidade: String
  fun ler(): Double
  fun formatado() = "${ler()} $unidade"  // padrão
}

// classes são final por padrão: use open
open class Sensor(val pino: Int) {
  open fun calibrar() { println("calibrando $pino") }
}

class Termistor(pino: Int) : Sensor(pino), Medidor {
  override val unidade = "°C"
  override fun ler() = 25.0
  override fun calibrar() {
    super.calibrar()
    println("ajustando curva beta")
  }
}
\end{codigo}
    \end{column}
    \begin{column}{0.45\textwidth}
      \centering
      \begin{tikzpicture}[node distance=7mm and 4mm]
        \node[caixaambar,minimum width=2.2cm] (m) {\textless\textless interface\textgreater\textgreater\\ \textbf{Medidor}};
        \node[caixa,minimum width=2.2cm,left=of m] (s) {\textbf{Sensor}\\ \tiny open};
        \node[caixaescura,minimum width=2.6cm,below=of $(s)!0.5!(m)$] (t) {\textbf{Termistor}};
        \draw[-{Triangle[open,length=2.5mm]},thick,petroleo] (t) -- (s);
        \draw[-{Triangle[open,length=2.5mm]},thick,dashed,ambar] (t) -- (m);
      \end{tikzpicture}
      \begin{conceito}[Regras]
        \begin{itemize}
          \item Classes e métodos são \textbf{final} por padrão; use \cod{open} para permitir herança.
          \item Herança simples de classe; múltiplas interfaces.
          \item Interfaces podem ter implementação padrão.
        \end{itemize}
      \end{conceito}
    \end{column}
  \end{columns}
\end{frame}

\begin{frame}[fragile]{\cod{data class}, \cod{object} e \cod{companion object}}
  \begin{columns}[T]
    \begin{column}{0.53\textwidth}
\begin{codigo}[title=Data class]{Kotlin}
data class Leitura(
  val sensorId: String,
  val valor: Double,
  val instante: Long = System.currentTimeMillis()
)

val a = Leitura("t1", 24.5, 1000)
val b = Leitura("t1", 24.5, 1000)
println(a == b)            // true: compara conteúdo
println(a)                 // Leitura(sensorId=t1, ...)
val c = a.copy(valor = 25.1)  // cópia alterada
val (id, valor) = a        // desestruturação
\end{codigo}
      \begin{conceito}[data class]
        Gera \cod{equals}, \cod{hashCode}, \cod{toString}, \cod{copy} e \cod{componentN}. É a base dos modelos e estados de tela.
      \end{conceito}
    \end{column}
    \begin{column}{0.45\textwidth}
\begin{codigo}[title=object e companion]{Kotlin}
// Singleton: uma única instância
object Config {
  const val BROKER = "mqtt.ifma.local"
  var debug = false
}

class Leitor private constructor(val id: String) {
  companion object {   // "estático"
    fun criar(id: String) = Leitor(id.trim())
  }
}

val l = Leitor.criar(" t1 ")
println(Config.BROKER)
\end{codigo}
    \end{column}
  \end{columns}
\end{frame}

\begin{frame}[fragile]{\cod{enum class} e \cod{sealed} para modelar estados}
  \begin{columns}[T]
    \begin{column}{0.47\textwidth}
\begin{codigo}[title=Enum]{Kotlin}
enum class StatusBancada(val cor: String) {
  LIVRE("verde"),
  RESERVADA("âmbar"),
  MANUTENCAO("vermelho");

  fun disponivel() = this == LIVRE
}

val s = StatusBancada.valueOf("LIVRE")
StatusBancada.entries.forEach { println(it) }
\end{codigo}
      \begin{dica}
        Use \cod{enum} quando os casos são fixos e \textbf{não carregam dados diferentes} entre si.
      \end{dica}
    \end{column}
    \begin{column}{0.51\textwidth}
\begin{codigo}[title=Sealed interface]{Kotlin}
sealed interface Resultado {
  data object Carregando : Resultado
  data class Sucesso(val dados: List<Leitura>)
    : Resultado
  data class Erro(val msg: String) : Resultado
}

// when exaustivo: sem else
fun render(r: Resultado) = when (r) {
  Resultado.Carregando -> "Carregando..."
  is Resultado.Sucesso -> "${r.dados.size} itens"
  is Resultado.Erro    -> "Falha: ${r.msg}"
}
\end{codigo}
      \begin{conceito}[Por que sealed?]
        Cada caso carrega dados próprios e o \cod{when} exaustivo obriga a tratar todos: um caso novo vira erro de compilação onde faltar tratamento.
      \end{conceito}
    \end{column}
  \end{columns}
\end{frame}

% =============================================================================
\section{Coleções}

\begin{frame}[fragile]{Listas, conjuntos e mapas}
  \begin{columns}[T]
    \begin{column}{0.52\textwidth}
\begin{codigo}[title=Coleções]{Kotlin}
// Somente leitura (preferível)
val nomes = listOf("DHT22", "BMP280", "LDR")
val ids = setOf(1, 2, 2, 3)          // {1, 2, 3}
val pinos = mapOf("led" to 2, "rele" to 4)

// Mutáveis
val fila = mutableListOf<Leitura>()
fila.add(Leitura("t1", 24.0))
fila += Leitura("t1", 24.2)

val cache = mutableMapOf<String, Double>()
cache["t1"] = 24.2
val v = cache["t2"] ?: 0.0     // get pode ser null

println(nomes[0])              // DHT22
println(nomes.getOrNull(9))    // null (sem exceção)
println(pinos.getValue("led")) // 2
\end{codigo}
    \end{column}
    \begin{column}{0.46\textwidth}
      \begin{conceito}[Leitura \emph{versus} mutável]
        \cod{List} não tem \cod{add}. \cod{MutableList} tem. Exponha sempre a versão de leitura e mantenha a mutável privada.
      \end{conceito}
      \footnotesize
      \begin{tabularx}{\linewidth}{@{}l L@{}}
        \toprule
        \cabecalho{Tipo} & \cabecalho{Característica} \\
        \midrule
        \cod{List} & ordenada, aceita repetidos \\
        \cod{Set} & sem repetidos \\
        \cod{Map} & pares chave $\rightarrow$ valor \\
        \bottomrule
      \end{tabularx}
    \end{column}
  \end{columns}
\end{frame}

\begin{frame}[fragile]{Operações funcionais}
  \begin{columns}[T]
    \begin{column}{0.56\textwidth}
\begin{codigo}[title=Pipeline de dados]{Kotlin}
val leituras = listOf(
  Leitura("t1", 24.1), Leitura("t2", 81.3),
  Leitura("t1", 24.9), Leitura("t3", 19.0),
  Leitura("t2", 79.8)
)

val alertas = leituras
  .filter { it.valor > 75 }               // seleciona
  .map { "${it.sensorId}: ${it.valor}" }  // transforma

val mediaPorSensor = leituras
  .groupBy { it.sensorId }
  .mapValues { (_, l) -> l.map { it.valor }.average() }

val maior = leituras.maxByOrNull { it.valor }
val total = leituras.sumOf { it.valor }
val desc = leituras.sortedByDescending { it.valor }
val temAlerta = leituras.any { it.valor > 80 }
val t3 = leituras.firstOrNull { it.sensorId == "t3" }
\end{codigo}
    \end{column}
    \begin{column}{0.42\textwidth}
      \scriptsize
      \begin{tabular}{@{}l l@{}}
        \toprule
        \cabecalho{Função} & \cabecalho{Resultado} \\
        \midrule
        \cod{map} & transforma cada item \\
        \cod{filter} & mantém os que passam \\
        \cod{groupBy} & agrupa em um \cod{Map} \\
        \cod{sortedBy} & nova lista ordenada \\
        \cod{sumOf}, \cod{average} & agregações \\
        \cod{any}, \cod{all}, \cod{none} & testes lógicos \\
        \cod{firstOrNull} & primeiro que satisfaz \\
        \cod{associateBy} & \cod{Map} por chave \\
        \cod{chunked}, \cod{windowed} & blocos e janelas \\
        \bottomrule
      \end{tabular}
      \begin{dica}
        \cod{windowed(5)} calcula médias móveis de leituras de sensor em uma linha.
      \end{dica}
    \end{column}
  \end{columns}
\end{frame}

% =============================================================================
\section{Funções de escopo e erros}

\begin{frame}[fragile]{Funções de escopo}
  \begin{columns}[T]
    \begin{column}{0.5\textwidth}
      \footnotesize
      \begin{tabularx}{\linewidth}{@{}l c c L@{}}
        \toprule
        \cabecalho{Função} & \cabecalho{Objeto} & \cabecalho{Retorna} & \cabecalho{Uso típico} \\
        \midrule
        \cod{let} & \cod{it} & lambda & tratar anuláveis \\
        \cod{apply} & \cod{this} & objeto & configurar objeto \\
        \cod{also} & \cod{it} & objeto & efeito extra (log) \\
        \cod{run} & \cod{this} & lambda & calcular com o objeto \\
        \cod{with} & \cod{this} & lambda & agrupar chamadas \\
        \bottomrule
      \end{tabularx}
    \end{column}
    \begin{column}{0.48\textwidth}
\begin{codigo}[title=Exemplos]{Kotlin}
val cfg = ConfigMqtt().apply {
  host = "broker.local"
  porta = 1883
}

val leitura = ler().also {
  Log.d(TAG, "lida: $it")
}

dispositivo?.let { conectar(it) }

val resumo = with(leitura) {
  "$sensorId = $valor"
}
\end{codigo}
    \end{column}
  \end{columns}
  \begin{dica}
    Use com moderação: muitas funções de escopo aninhadas prejudicam a leitura.
  \end{dica}
\end{frame}

\begin{frame}[fragile]{Tratamento de erros}
  \begin{columns}[T]
    \begin{column}{0.5\textwidth}
\begin{codigo}[title=try é expressão]{Kotlin}
val valor: Double = try {
  entrada.toDouble()
} catch (e: NumberFormatException) {
  Log.w(TAG, "valor inválido: $entrada")
  0.0
} finally {
  fecharArquivo()
}

fun lerSensor(): Double {
  if (!conectado) throw IllegalStateException("desconectado")
  return 24.5
}
\end{codigo}
    \end{column}
    \begin{column}{0.48\textwidth}
\begin{codigo}[title=Result e runCatching]{Kotlin}
val r: Result<Double> = runCatching { lerSensor() }

r.onSuccess { println("ok: $it") }
 .onFailure { println("falhou: ${it.message}") }

val v = r.getOrDefault(Double.NaN)
val txt = r.fold(
  onSuccess = { "%.1f".format(it) },
  onFailure = { "--" }
)
\end{codigo}
      \begin{conceito}
        Kotlin não tem \emph{checked exceptions}: o compilador não obriga a capturar nada. Por isso é importante documentar e tratar falhas de rede, arquivo e conversão.
      \end{conceito}
    \end{column}
  \end{columns}
\end{frame}

% =============================================================================
\section{Fechamento}

\begin{frame}{Resumo}
  \begin{resumo}
    \begin{itemize}
      \item Prefira \cod{val}; a inferência de tipos deixa o código enxuto sem perder segurança.
      \item \cod{T?} marca valores anuláveis; use \cod{?.}, \cod{?:}, \cod{let} e \emph{smart casts}, e evite \cod{!!}.
      \item \cod{if}, \cod{when} e \cod{try} são expressões que devolvem valor.
      \item Parâmetros padrão e nomeados, extensões e lambdas são a base da API do Compose.
      \item \cod{data class} para dados, \cod{sealed} para estados, \cod{object} para singletons.
      \item Coleções somente leitura + \cod{map}/\cod{filter}/\cod{groupBy} substituem a maioria dos laços.
    \end{itemize}
  \end{resumo}
\end{frame}

\begin{frame}{Glossário}
  \begin{glossario}
    Paradigma & estilo de organizar o código (imperativo, objetos, funcional, declarativo) \\
    Tipagem estática & tipos verificados na compilação \\
    Inferência de tipos & o compilador deduz o tipo a partir do valor \\
    \emph{Null safety} & distinção entre tipos anuláveis (\cod{String?}) e não anuláveis \\
    \emph{Smart cast} & conversão automática após uma verificação de tipo ou de nulo \\
    \cod{data class} & classe de dados com \cod{equals}, \cod{hashCode}, \cod{copy} e \cod{toString} gerados \\
    \cod{sealed class} & hierarquia fechada: o \cod{when} sabe todos os casos \\
    Função de ordem superior & função que recebe ou devolve outra função \\
    Coleta de lixo & liberação automática de objetos sem referências \\
  \end{glossario}
\end{frame}

\begin{frame}{Questões de revisão}
  \begin{questoes}
    \item Que vantagens a tipagem estática traz para um app que vai crescer?
    \item Por que o nulo foi chamado de ``erro de um bilhão de dólares''? Como o Kotlin o trata?
    \item Qual a diferença entre \cod{val} e \cod{var}, e por que preferir \cod{val}?
    \item Dê um exemplo de encapsulamento e um de polimorfismo no LabMonitor.
    \item O que é uma função pura? Por que ela é fácil de testar?
    \item Para que serve uma \cod{sealed class} ao modelar o estado de uma tela?
    \item Explique a diferença entre igualdade de valor e de identidade com \cod{data class}.
    \item O que acontece com um objeto quando nenhuma variável aponta mais para ele?
  \end{questoes}
\end{frame}

\begin{frame}{Exercícios}
  \begin{exercicio}[Exercício 1 --- Conversões]
    Crie funções de extensão \cod{Double.cParaF()}, \cod{Double.fParaC()} e \cod{Int.adcParaVolts(vref: Double = 3.3, bits: Int = 12)}, que converte uma leitura de ADC em tensão.
  \end{exercicio}
  \begin{exercicio}[Exercício 2 --- Coleções]
    Dada uma lista de \cod{Leitura}, calcule: a média por sensor; o sensor com maior variância; a média móvel de 5 amostras do sensor \cod{"t1"}; as leituras fora da faixa 10--40\,°C.
  \end{exercicio}
  \begin{exercicio}[Exercício 3 --- Modelagem]
    Modele com \cod{sealed interface} os estados de uma conexão BLE: \emph{Desconectado}, \emph{Procurando}, \emph{Conectando(endereco)}, \emph{Conectado(endereco, rssi)} e \emph{Falha(motivo)}. Escreva uma função que gere a mensagem de cada estado com um \cod{when} exaustivo.
  \end{exercicio}
  \begin{projeto}[LabMonitor --- Etapa 2]
    Crie o pacote \cod{model} com \cod{Equipamento}, \cod{Bancada}, \cod{Leitura} e \cod{enum StatusBancada}, mais uma lista de dados fictícios em \cod{object DadosExemplo}.
  \end{projeto}
\end{frame}

\begin{frame}{Para aprofundar}
  \small
  \begin{itemize}
    \item Documentação oficial: \url{https://kotlinlang.org/docs/home.html}
    \item Kotlin Koans: \url{https://play.kotlinlang.org/koans}
    \item Kotlin para desenvolvedores Android: \url{https://developer.android.com/kotlin}
    \item Convenções de código: \url{https://kotlinlang.org/docs/coding-conventions.html}
    \item Livro gratuito: \emph{Kotlin Bootcamp for Programmers} (codelab do Google)
  \end{itemize}
\end{frame}

\begin{frame}{Respostas das questões de revisão}
  \begin{respostas}
    \item Erros de tipo aparecem na compilação, a IDE completa e refatora com segurança e os tipos documentam a intenção do código.
    \item Porque referências nulas causaram incontáveis falhas e custos. O Kotlin separa \cod{String} de \cod{String?} e exige tratar o nulo (\cod{?.}, \cod{?:}, \emph{smart cast}) antes do uso.
    \item \cod{val} não pode ser reatribuída; \cod{var} pode. \cod{val} evita mudanças inesperadas e facilita o raciocínio sobre o código.
    \item Encapsulamento: \cod{Reserva} só muda por \cod{cancelar()}. Polimorfismo: \cod{FonteBle} e \cod{FonteFake} respondem a \cod{leituras()} cada uma do seu jeito.
    \item Devolve sempre o mesmo resultado para a mesma entrada, sem efeitos colaterais; o teste só precisa conferir entrada e saída.
    \item Lista todos os estados possíveis (carregando, sucesso, erro); o \cod{when} exaustivo obriga a tratar cada um.
    \item Duas instâncias de \cod{Bancada("B3")} são iguais por valor (\cod{==} verdadeiro), mas não são o mesmo objeto (\cod{===} falso).
    \item Ele se torna inalcançável e é liberado pelo coletor de lixo em algum momento.
  \end{respostas}
\end{frame}

\encerramento{Módulo 03 (comum) --- Engenharia de software e Git: requisitos, métodos ágeis, versionamento e colaboração.}

\end{document}
